% Responsável: P1 + P2. Enunciado: >= 8 referências, COMENTADAS CRITICAMENTE (não só enumeradas).
% Factos verificados nos resumos/textos dos artigos em 17/11/2026 (Potthast, Pérez-Rosas, Horne & Adalı,
% Santos et al.). Ler os artigos completos antes da apresentação.
\section{Trabalho relacionado}\label{sec:rel}

\paragraph{Panoramas.}
Shu et~al.~\cite{shu2017} e Zhou e Zafarani~\cite{zhou2020} organizam a deteção de notícias falsas
em abordagens baseadas no conteúdo (conhecimento e estilo), na propagação e na fonte. O nosso
trabalho situa-se na deteção pelo \emph{estilo} do conteúdo, a única possível quando, como no
cenário da organização, só se dispõe do texto. A nossa análise de erros confirma a fragilidade que
estes panoramas apontam a esta família de métodos: basta imitar a escrita jornalística para enganar
o detetor.

\paragraph{Estilo e notícias falsas.}
Horne e Adal{\i}~\cite{horne2017} concluíram que, em inglês, as notícias falsas se parecem mais com
a sátira do que com notícias reais, que a estrutura do título e o uso de nomes próprios as distinguem
e que persuadem por heurísticas, dirigindo-se a leitores que não passam do título. Potthast
et~al.~\cite{potthast2018} mostraram que o estilo separa bem notícias hiperpartidárias das restantes
(F1 de 0,78) e a sátira de ambas (0,81), mas que a deteção de notícias falsas apenas pelo estilo
falha (F1 de 0,46). A diferença face aos valores próximos de 0,95 no Fake.Br é instrutiva: quando as
notícias falsas vêm de poucos sítios, o estilo da fonte confunde-se com a falsidade. Pérez-Rosas
et~al.~\cite{perezrosas2018} construíram dois conjuntos em inglês (FakeNewsAMT e Celebrity) e
combinaram atributos lexicais, sintáticos e de legibilidade; ao treinar num domínio e testar no
outro, a exatidão desceu de 0,67--0,78 para 0,51--0,61. No nosso protocolo por tema a queda é muito
menor (até 5,5 pontos), possivelmente porque as fontes são as mesmas em todos os temas.

\paragraph{Notícias falsas em português.}
Monteiro et~al.~\cite{monteiro2018} criaram o Fake.Br: 3\,600 notícias falsas e 3\,600 verdadeiras,
alinhadas por assunto e recolhidas entre janeiro de 2016 e janeiro de 2018, com 89\% de exatidão
usando SVM sobre \emph{bag-of-words}; atributos como a pausalidade e a emotividade, sozinhos, ficaram
perto de 55\%~\cite{santos2020}. Silva et~al.~\cite{silva2020} disponibilizaram o conjunto
publicamente e compararam vários métodos de aprendizagem automática, conjuntos de atributos e
classificadores. Santos et~al.~\cite{santos2020} mostraram que atributos de legibilidade, sozinhos,
chegam a 92\% de exatidão, com SVM e validação cruzada de 5 folds. Estes trabalhos estabeleceram o
corpus e as referências de desempenho, mas não indicam se as duas notícias de cada par ficam no mesmo
fold, nem quanto do desempenho se explica pelo comprimento ou pela forma como os textos acabam. São
estas as questões que o nosso protocolo controla: dividimos por par, medimos as \emph{baselines} de
comprimento e de final do texto e testamos temas não vistos.

\paragraph{Representações e avaliação.}
Usamos TF-IDF~\cite{manning2008} e o BERTimbau~\cite{souza2020}, uma variante do BERT~\cite{devlin2019}
pré-treinada em português do Brasil. Para os tópicos usamos LDA~\cite{blei2003}, avaliado com a
coerência NPMI, que Röder et~al.~\cite{roder2015} identificaram como uma das medidas mais
correlacionadas com o julgamento humano. Kapoor e Narayanan~\cite{kapoor2023} documentam como a
fuga de informação entre treino e teste inflaciona resultados em muitas áreas; no Fake.Br, o
emparelhamento das notícias torna esse risco concreto.
